\usepackage[utf8]{inputenc}
\usepackage[T1]{fontenc}
\usepackage[spanish]{babel}

% Tamaño de letra 12pt (en documentclass) y espaciado 1,5
\usepackage{setspace}
\onehalfspacing

\usepackage{graphicx}
\usepackage{float} % [H] para fijar figuras y que no floten a otras secciones
\usepackage{placeins} % \FloatBarrier para evitar que flotantes pasen a la siguiente sección
\usepackage{eso-pic}
\usepackage{changepage} % Para ajustar márgenes
\usepackage[margin=2.5cm]{geometry} % Márgenes de 2.5 cm en todos los lados

% Indentación completa de párrafos en subsecciones (todas las líneas, no solo la primera)
\newenvironment{indentedsubsection}{%
  \begin{adjustwidth}{1cm}{0cm}%
  \setlength{\parindent}{0pt}%
}{%
  \end{adjustwidth}%
}
\usepackage{array}
\usepackage{tabularx}

% Hipervínculos (índice clicable)
\usepackage[hidelinks]{hyperref}

% Títulos y subtítulos del mismo tamaño que el texto
\usepackage{titlesec}
\titleformat{\section}{\normalsize\bfseries}{\thesection. }{0.5em}{}
\titleformat{\subsection}{\normalsize\bfseries}{\thesubsection. }{0.5em}{}

% Nombres en español para listas de figuras y tablas
\renewcommand{\listfigurename}{Índice de ilustraciones}
\renewcommand{\listtablename}{Índice de tablas}
